\documentclass[11pt]{article}

\input{../shared/preamble.tex}
\usepackage[margin=1.15in]{geometry}
\usepackage{enumitem}
\usepackage{xcolor}
\usepackage{tcolorbox}
\setlist{nosep, leftmargin=*}

\title{Hierarchical Core-Orbital Allocation\\[2pt]
\large What Can Be Recovered from Public Sources}
\author{Notes compiled by Peter Cotton}
\date{September 16, 2026}

\begin{document}
\maketitle

\begin{tcolorbox}[colback=gray!6, colframe=gray!50, boxrule=0.4pt, arc=2pt]
\small
\textbf{The paper.} Mario Bajo Traver, \emph{Hierarchical Core-Orbital
Allocation: Decoupling Clustering from Optimization in Portfolio
Construction}. Forthcoming in \emph{The Journal of Financial Data Science}
(Portfolio Management Research). DOI 10.2139/ssrn.5928122.

\medskip
\textbf{The author.} Head of Investment Strategies Unit, Department of
Market Analysis and Intelligence, Banco de Espa\~na. No personal homepage.
ORCID 0009-0006-3791-2216 points back to SSRN.

\medskip
\textbf{Provenance.} The full text is not publicly fetchable as of the date
above. SSRN hosts the abstract sheet only. ResearchGate publication 398769353
(updated 19~December~2025) holds the full text but blocks automated access.
Everything below comes from the abstract, the author's LinkedIn
announcement (\url{https://lnkd.in/p/ekyaNa-E}), and search-engine excerpts
of the ResearchGate full text. Passages in quotation marks are verbatim
from those sources.
\end{tcolorbox}

\section*{Architecture}

Two levels.

\paragraph{Strategic.} Cluster the universe. Select one representative
``core'' asset per cluster. Solve the allocation ``operating exclusively on
core assets, which constitute a representative set of reduced dimension'',
with any objective (minimum variance, maximum Sharpe, risk parity), ``a
formal optimization or a closed-form allocation''.

\paragraph{Tactical.} Strategic weights ``are redistributed to individual
assets within each cluster (`orbitals') through hierarchical rules or local
optimization, preserving global coherence: the sum of intra-cluster weights
exactly reproduces the upper-level allocation.''

\section*{Core selection}

Three criteria are described.

\begin{itemize}
\item \emph{Medoid.} ``Unlike the centroid (geometric center), the medoid is
  always an element of the original set, conferring robustness against
  outliers''; it ``captures the central position in correlation space''.
\item \emph{Principal-component alignment.} ``Selects the asset whose
  dynamics are most aligned with the cluster's first principal component.''
  ``High values (PCA $> 0.7$) indicate that joint variation is predominantly
  explained by a single factor.''
\item \emph{CS$^2$, explanatory capacity.} ``Selects the asset that maximizes
  the sum of squared correlations with other cluster members.''
\end{itemize}

Guidance given: ``in heterogeneous clusters, medoid is robust to subgroups
or nonlinear structures, PCA is optimal when a clear dominant factor
exists, and CS$^2$ maximizes internal predictive power.''

Once the $k$ cores are fixed, ``the set of orbital assets for cluster $c$''
is the remainder of the cluster.

\section*{Claimed properties}

\begin{enumerate}
\item \emph{Numerical robustness.} ``Operating on a reduced covariance space
  ($k$ clusters $< N$ assets) with improved condition number.''
\item \emph{Functional flexibility.} ``Any combination of objectives at both
  levels without altering the structural logic.''
\item \emph{Hierarchical coherence.} ``Intra-cluster decisions respect the
  upper-level allocation'', enforced ``through aggregation constraints''.
\end{enumerate}

Stated contrast with HRP and HERC: they ``optimize directly based on
dendrogram topology, implementing an allocation logic that cannot be
modified without altering the entire algorithm.''

The clustering is described as ``stable in dimensionality yet flexible in
allocation'', with ``a practically stable number of clusters'' over time.

\section*{Empirical setup}

\begin{center}
\begin{tabular}{@{}ll@{}}
\toprule
Universe & USD-denominated bonds: sovereign, supranational, corporate \\
         & ``multiple maturities, issuers and credit ratings'' \\
Sample   & July 2014 to September 2025, 135 months \\
Rebalancing & 45 rebalancings (quarterly), out of sample \\
Baselines & NCO, HRP, HERC, mean--variance, equal weight \\
Profiles & conservative, balanced, aggressive \\
\bottomrule
\end{tabular}
\end{center}

``Assets naturally group by issuer type, credit quality, and duration.''

Headline result: ``risk-adjusted performance is statistically comparable to
NCO, the closest two-stage alternative, while HCOA delivers greater
diversification and, on average, lower constraint-projection distance.
Results are robust to alternative clustering parameter choices.''

\section*{Not recovered}

Clustering distance and linkage, how $k$ is chosen, the exact tactical
formulas, and the numerical tables.

A third-party replication card at \texttt{paperswithbacktest.com} (48
bonds, 1990 to 2026, quarterly) reports a Sharpe ratio of $0.25$ and a
maximum drawdown of $8.1\%$. Its configuration is unknown and it is not
evidence about the paper.

\end{document}
